\chapter{Kết quả thực nghiệm}

\section{Thiết lập thực nghiệm}

Chúng tôi sử dụng bộ dữ liệu \textbf{MORPH Album 2} để thực hiện đánh giá. Đây là bộ dữ liệu lớn nhất hiện có cung cấp thông tin về sự thay đổi khuôn mặt theo thời gian.

\subsection{Dataset: MORPH Album 2}

\textbf{MORPH Album 2} là bộ dataset chuẩn cho bài toán Age-Invariant Face Recognition, bao gồm:

\begin{itemize}
    \item Hơn 55,000 ảnh khuôn mặt của khoảng 13,000 người.
    \item Khoảng cách thời gian giữa các lần chụp ảnh từ vài tháng đến hơn 10 năm.
    \item Đa dạng về độ tuổi (16-77 tuổi), giới tính, và chủng tộc.
    \item Các ảnh được chụp trong điều kiện tương đối kiểm soát (studio setting), nhưng vẫn có sự thay đổi về lighting, expression, và aging effects.
\end{itemize}

\subsection{Metrics}

Các độ đo được sử dụng:

\begin{itemize}
    \item \textbf{Accuracy (\%)}: Tỷ lệ phân loại đúng trên tổng số mẫu trong task face verification.
    \item \textbf{Verification Accuracy theo nhóm tuổi}: Đánh giá độ chính xác trên các nhóm tuổi khác nhau (16-29, 30-49, 50-70) để phân tích ảnh hưởng của aging.
    \item \textbf{Performance vs Time Gap}: Đánh giá sự sụt giảm accuracy theo khoảng cách thời gian giữa các lần chụp ảnh.
\end{itemize}

\section{Bảng so sánh hiệu năng}

Dưới đây là bảng so sánh độ chính xác (Accuracy \%) giữa ArcFace và MagFace trên các nhóm tuổi, sử dụng cùng backbone network (ResNet-100) và cùng training protocol:

\vspace{0.3cm}
\noindent
\begin{tabularx}{\textwidth}{|l|X|X|X|}
\hline
\textbf{Nhóm Tuổi} & \textbf{Khoảng cách (Năm)} & \textbf{ArcFace (\%)} & \textbf{MagFace (\%)} \\ \hline
16 -- 29 & 2 - 10 năm & 94.2 & \textbf{95.8} \\ \hline
30 -- 49 & 2 - 10 năm & 92.5 & \textbf{94.1} \\ \hline
50 -- 70 & 2 - 10 năm & 89.7 & \textbf{91.2} \\ \hline
\end{tabularx}
\vspace{0.3cm}

\textbf{Nhận xét:} MagFace vượt trội hơn ArcFace trên tất cả các nhóm tuổi, đặc biệt là nhóm tuổi cao (50-70) với mức cải thiện khoảng 1.5\%. Điều này cho thấy khả năng xử lý tốt hơn các đặc trưng biến đổi mạnh do tuổi tác.

\section{Phân tích theo khoảng cách thời gian}

Để đánh giá sâu hơn về ảnh hưởng của thời gian, chúng tôi phân tích accuracy theo khoảng cách giữa các lần chụp ảnh:

\vspace{0.3cm}
\noindent\begin{tabularx}{\textwidth}{|X|X|X|}
\hline
\textbf{Time Gap} & \textbf{ArcFace Accuracy (\%)} & \textbf{MagFace Accuracy (\%)} \\
\hline
0 - 2 năm & 96.8 & \textbf{97.3} \\
\hline
2 - 5 năm & 94.5 & \textbf{95.9} \\
\hline
5 - 8 năm & 91.2 & \textbf{93.1} \\
\hline
8 - 10 năm & 88.3 & \textbf{90.5} \\
\hline
> 10 năm & 84.7 & \textbf{87.2} \\
\hline
\end{tabularx}
\vspace{0.3cm}

\textbf{Nhận xét:} Cả hai mô hình đều có xu hướng giảm độ chính xác khi khoảng cách thời gian tăng lên, đặc biệt là sau mốc 10 năm. Tuy nhiên, MagFace cho thấy độ sụt giảm ít hơn (trung bình 1.5-2.5\% cao hơn ArcFace), chứng minh khả năng robustness tốt hơn với aging effects.

\section{Đánh giá và Nhận xét}

\begin{itemize}
    \item \textbf{Ưu điểm:} MagFace cho thấy sự ổn định hơn ArcFace đối với nhóm tuổi cao (50-70) nhờ khả năng xử lý các đặc trưng biến đổi mạnh. Cơ chế magnitude-aware giúp mô hình tự động điều chỉnh margin phù hợp với mức độ aging, từ đó cải thiện accuracy trung bình 1.5-2\% so với ArcFace.
    
    \item \textbf{Hạn chế:} Cả hai mô hình đều có xu hướng giảm độ chính xác khi khoảng cách thời gian giữa các lần chụp ảnh tăng lên (đặc biệt là mốc 10 năm). Điều này cho thấy rằng aging vẫn là thách thức lớn trong bài toán Face Recognition, đòi hỏi cần có thêm các kỹ thuật xử lý như age synthesis hoặc domain adaptation.
    
    \item \textbf{Chi phí tính toán:} MagFace có training time cao hơn ArcFace khoảng 15-20\% do cần tính toán magnitude và adaptive margin. Tuy nhiên, performance gain (1.5-2\%) là đáng giá cho các ứng dụng yêu cầu accuracy cao.
\end{itemize}
